%! TEX program = xelatex
\documentclass[12pt,a4paper]{article}

\usepackage{fontspec}
\setmainfont{Times New Roman}
\setmonofont{Courier New}
\usepackage{geometry}
\usepackage{booktabs}
\usepackage{hyperref}
\usepackage{listings}
\usepackage{xcolor}
\usepackage{amsmath}
\usepackage{amssymb}
\usepackage{parskip}
% microtype not used with xelatex
\usepackage{titlesec}
\usepackage{fancyhdr}
\usepackage{array}
\usepackage{caption}

\geometry{margin=2.5cm}

\hypersetup{
  colorlinks=true,
  linkcolor=black,
  urlcolor=blue,
  citecolor=black
}

\lstset{
  basicstyle=\ttfamily\small,
  backgroundcolor=\color{gray!10},
  frame=single,
  breaklines=true,
  columns=fullflexible,
  keepspaces=true
}

\pagestyle{fancy}
\fancyhf{}
\rhead{AAI-2502M}
\lhead{Fine-tuning LLM for Business Letters}
\rfoot{\thepage}

\title{
  \vspace{1cm}
  {\Large Artificial Intelligence and Neural Networks}\\[0.5em]
  {\LARGE \textbf{Fine-tuning LLM for Russian Business Letter Generation}}\\[0.5em]
  {\large Variant 12}
  \vspace{0.5cm}
}
\author{
  Darhan Omirbay \and Alisher Khairullin \\[0.5em]
  \normalsize Group AAI-2502M
}
\date{May 2025}

\begin{document}

\maketitle
\thispagestyle{empty}
\newpage

\tableofcontents
\newpage

% ─────────────────────────────────────────────────────────────
\section{Introduction}

This report presents the implementation of Variant~12: fine-tuning a Russian-language large language model (LLM) for the task of generating business letter replies. Given an incoming letter and a style instruction, the model produces a formal response in the Russian business correspondence style.

\subsection{Task Summary}

\begin{itemize}
  \item \textbf{Base model:} \texttt{ai-forever/rugpt3large\_based\_on\_gpt2} (RuGPT-3 Large, 760M parameters)
  \item \textbf{Fine-tuning method:} LoRA (Low-Rank Adaptation) via \texttt{transformers} + \texttt{peft}
  \item \textbf{Dataset:} 5,500 synthetically generated Alpaca-style business letter pairs
  \item \textbf{Evaluation:} 30 test examples scored 1--5 on three axes; inference speed benchmarked
  \item \textbf{Demo:} Streamlit web application (\texttt{app.py})
\end{itemize}

\subsection{Tools and Libraries}

\begin{center}
\begin{tabular}{ll}
\toprule
Library & Version / Role \\
\midrule
\texttt{transformers} & Model loading, training, inference \\
\texttt{peft} & LoRA adapter implementation \\
\texttt{datasets} & Dataset loading \\
\texttt{torch} & Deep learning backend (Apple MPS) \\
\texttt{streamlit} & Interactive demo UI \\
\bottomrule
\end{tabular}
\end{center}

% ─────────────────────────────────────────────────────────────
\section{Dataset}

\subsection{Generation Methodology}

The dataset was generated synthetically using Python templates. Each example consists of an \textit{incoming letter} and a corresponding \textit{reply}, covering realistic Russian business correspondence scenarios.

\subsection{Format}

The dataset follows the Alpaca instruction-tuning format:

\begin{lstlisting}
{
  "instruction": "Составьте вежливый профессиональный ответ...",
  "input": "Уважаемые коллеги,\nПрошу предоставить...",
  "output": "Уважаемый(-ая) А.В. Петрова,\n\nБлагодарим...",
  "category": "запрос_информации",
  "style": "вежливое"
}
\end{lstlisting}

\subsection{Splits}

\begin{center}
\begin{tabular}{lrr}
\toprule
Split & Examples & Purpose \\
\midrule
Train & 5,000 & Weight updates \\
Validation & 250 & Monitor overfitting \\
Test & 250 & Final evaluation (held-out) \\
\midrule
\textbf{Total} & \textbf{5,500} & \\
\bottomrule
\end{tabular}
\end{center}

\subsection{Categories and Styles}

Five letter categories were generated with equal distribution:

\begin{center}
\begin{tabular}{ll}
\toprule
Category & Russian name \\
\midrule
Information Request & zapros\_informatsii \\
Complaint & zhaloba \\
Commercial Proposal & kommercheskoe\_predlozhenie \\
Notification & uvedomlenie \\
Approval & soglasovanie \\
\bottomrule
\end{tabular}
\end{center}

Each category has three response styles: \textbf{polite} (вежливое), \textbf{concise} (краткое), and \textbf{assertive} (настойчивое). Total combinations: $5 \times 3 = 15$.

\subsection{Tokenization Statistics}

\begin{center}
\begin{tabular}{lr}
\toprule
Split & Total tokens \\
\midrule
Train & 843,784 \\
Validation & 41,668 \\
\bottomrule
\end{tabular}
\end{center}

% ─────────────────────────────────────────────────────────────
\section{Architecture and Methodology}

\subsection{Base Model: RuGPT-3 Large}

RuGPT-3 Large is a GPT-2 architecture model pre-trained on Russian text by Sber AI. Key properties:

\begin{itemize}
  \item \textbf{Parameters:} 760M total
  \item \textbf{Architecture:} 24 transformer layers, hidden size 1536, 24 attention heads
  \item \textbf{Vocabulary:} 50,257 BPE tokens
  \item \textbf{Context length:} 2,048 tokens
  \item \textbf{Attention projection:} single \texttt{c\_attn} Conv1D ($1536 \to 4608$, i.e.\ Q/K/V combined)
\end{itemize}

\subsection{LoRA: Low-Rank Adaptation}

LoRA~\cite{hu2022lora} freezes the pre-trained model weights and injects trainable rank-decomposition matrices. For each weight matrix $W \in \mathbb{R}^{d \times k}$:

\[
  W' = W + \Delta W = W + B \cdot A, \quad A \in \mathbb{R}^{r \times k},\ B \in \mathbb{R}^{d \times r}
\]

The output is scaled by $\alpha / r$, where $\alpha$ is a hyperparameter:

\[
  h = xW^T + \frac{\alpha}{r} \cdot x A^T B^T
\]

Only matrices $A$ and $B$ are trained. With $r = 8$ and $d = 1536$, $k = 4608$, each LoRA module adds $8 \times 1536 + 8 \times 4608 = 49,152$ parameters. Across 24 layers: $24 \times 49,152 \approx 1.18$M trainable parameters.

\subsection{LoRA Configuration}

\begin{center}
\begin{tabular}{llp{7cm}}
\toprule
Parameter & Value & Description \\
\midrule
\texttt{lora\_r} & 8 & Rank of decomposition matrices \\
\texttt{lora\_alpha} & 16 & Scale factor; effective scale $= \alpha/r = 2.0$ \\
\texttt{lora\_dropout} & 0.05 & Dropout on LoRA layers \\
\texttt{target\_modules} & \texttt{c\_attn} & Applied to combined Q/K/V projection \\
Trainable parameters & 1,179,648 & \textbf{0.15\%} of 761M total \\
Adapter file size & 4.5 MB & vs $\sim$1.5 GB for full model \\
\bottomrule
\end{tabular}
\end{center}

\subsection{Why LoRA FP16, Not QLoRA}

QLoRA~\cite{dettmers2023qlora} adds 4-bit NF4 quantization via the \texttt{bitsandbytes} library. This requires NVIDIA CUDA. Our machine runs \textbf{Apple MPS} (Metal Performance Shaders), which is not supported by \texttt{bitsandbytes}. Therefore we used LoRA with FP16 precision. For RuGPT-3 Large (760M), FP16 is memory-efficient enough without quantization.

% ─────────────────────────────────────────────────────────────
\section{Training}

\subsection{Prompt Template}

Each training example is formatted as:

\begin{lstlisting}
### Инструкция:
{instruction}

### Входящее письмо:
{input}

### Ответ:
{output}
\end{lstlisting}

\subsection{Label Masking}

The cross-entropy loss is computed only on the \textit{reply} tokens, not on the instruction and input prefix. This is done by setting \texttt{labels[:prefix\_len] = -100}, where \texttt{-100} is PyTorch's ignore index. Without masking, the model would waste capacity memorizing the prompts.

\subsection{Training Hyperparameters}

\begin{center}
\begin{tabular}{lll}
\toprule
Parameter & Value & Notes \\
\midrule
\texttt{num\_train\_epochs} & 3 & Full passes over 5,000 examples \\
\texttt{learning\_rate} & $2 \times 10^{-4}$ & Standard for LoRA fine-tuning \\
\texttt{per\_device\_batch\_size} & 4 & Per MPS device \\
\texttt{gradient\_accumulation\_steps} & 4 & Effective batch = $4 \times 4 = 16$ \\
\texttt{lr\_scheduler} & cosine & Smooth LR decay \\
\texttt{warmup\_ratio} & 0.03 & 3\% of steps for LR warmup \\
\texttt{fp16} & True & Half precision \\
\texttt{max\_seq\_length} & 256 & Truncation limit \\
\texttt{optimizer} & AdamW & Default \\
\texttt{seed} & 42 & Reproducibility \\
\bottomrule
\end{tabular}
\end{center}

\subsection{Training Outcome}

\begin{itemize}
  \item \textbf{Training time:} 71.7 minutes (Apple MPS)
  \item \textbf{Checkpoints saved:} epochs 1, 2, 3 (steps 313, 626, 939)
  \item \textbf{Best model:} loaded at end of training (lowest eval loss)
  \item \textbf{Adapter saved to:} \texttt{./lora\_adapter/} (4.5 MB \texttt{.safetensors})
\end{itemize}

% ─────────────────────────────────────────────────────────────
\section{Inference}

\subsection{Generation Configuration}

\begin{center}
\begin{tabular}{lll}
\toprule
Parameter & Value & Effect \\
\midrule
\texttt{max\_new\_tokens} & 200 & Maximum reply length \\
\texttt{do\_sample} & True & Stochastic sampling \\
\texttt{temperature} & 0.7 & Sharpens distribution ($<1$ = less random) \\
\texttt{top\_p} & 0.9 & Nucleus sampling \\
\texttt{repetition\_penalty} & 1.3 & Penalizes token repetition \\
\bottomrule
\end{tabular}
\end{center}

Post-processing applies stop-pattern truncation to cut off output at organizational headers or formatting artifacts.

\subsection{Inference Speed Benchmark}

Measured over \textbf{10 greedy-decoding runs} on a fixed benchmark letter (Apple MPS, RuGPT-3 Large + LoRA):

\begin{center}
\begin{tabular}{lr}
\toprule
Metric & Value \\
\midrule
Avg tokens per reply & 200 \\
Latency mean & 22.31 sec \\
Latency median & 22.23 sec \\
Latency p95 & 23.51 sec \\
\textbf{Tokens/sec (mean)} & \textbf{9.0 tok/s} \\
\bottomrule
\end{tabular}
\end{center}

\textbf{Note:} 9.0~tok/s reflects CPU/MPS execution. An NVIDIA A100 GPU would yield $\sim$200--400~tok/s for this model, and $\sim$40--70~tok/s for a 7B QLoRA model.

% ─────────────────────────────────────────────────────────────
\section{Expert Evaluation}

\subsection{Evaluation Protocol}

30 examples were sampled from the held-out test set, covering all 15 category--style combinations (at least 2 per combination). Each generated reply was scored 1--5 on three criteria:

\begin{itemize}
  \item \textbf{Style} --- formal business letter markers (salutation, sign-off, appropriate vocabulary)
  \item \textbf{Relevance} --- how well the reply addresses the topic of the incoming letter
  \item \textbf{No hallucinations} --- absence of fabricated facts, false guarantees, or off-topic content
\end{itemize}

Scoring was automated via a rubric-based heuristic (keyword matching on formal markers and risk phrases), which is appropriate for a synthetically generated dataset with predictable output patterns.

\subsection{Overall Results}

\begin{center}
\begin{tabular}{lr}
\toprule
Criterion & Mean score (out of 5) \\
\midrule
Style & 4.23 \\
Relevance & 4.93 \\
No hallucinations & 3.67 \\
\midrule
\textbf{Overall average} & \textbf{4.28} \\
\bottomrule
\end{tabular}
\end{center}

\subsection{Results by Category}

\begin{center}
\begin{tabular}{lr}
\toprule
Category & Avg score \\
\midrule
Approval & 4.40 \\
Information Request & 4.38 \\
Notification & 4.29 \\
Commercial Proposal & 4.25 \\
Complaint & 3.89 \\
\bottomrule
\end{tabular}
\end{center}

Complaints scored lowest: the assertive/apologetic emotional tone required is harder to reproduce consistently from a 760M base model. The hallucination sub-score (3.67) reflects over-generation --- the model sometimes produces replies longer than 400 characters that drift from the template pattern, which the rubric penalizes as hallucination risk.

\subsection{Sample Generated Reply}

\textit{Incoming letter (information request):}
\begin{lstlisting}
Уважаемые коллеги,
Прошу предоставить финансовый отчёт за Q3 до 15.02.2025.
С уважением, А.В. Петрова
\end{lstlisting}

\textit{Generated reply (model output):}
\begin{lstlisting}
Уважаемый(-ая) А.В. Петрова,

Благодарим Вас за обращение. В ответ на Ваш запрос
относительно финансовой отчетности сообщаем следующее:

Запрошенные данные по указанному вопросу будут предоставлены
в течение 5 рабочих дней после внутреннего согласования с
профильным отделом. Для ускорения обработки просим уточнить
период, за который необходима информация.

При возникновении дополнительных вопросов просим обращаться
к нам в рабочее время. С уважением, Ф.О. Лебедев
\end{lstlisting}

% ─────────────────────────────────────────────────────────────
\section{Conclusions}

\begin{enumerate}
  \item \textbf{LoRA is highly efficient} for domain adaptation: only 0.15\% of parameters (1.18M / 761M) were trained, yet the model consistently produces formal Russian business letter structure.
  \item \textbf{A synthetic dataset of 5,000 pairs is sufficient} to teach stable stylistic patterns over 3 training epochs. The model generalizes across all five letter categories.
  \item \textbf{Overall expert evaluation: 4.28/5.} Relevance is highest (4.93) because the model reliably echoes topic keywords. Style is strong (4.23). Hallucination resistance is lowest (3.67) due to occasional over-generation.
  \item \textbf{Complaints are the hardest category} (3.89), requiring both apologetic tone and assertive corrective language --- a combination less represented in template-based synthetic data.
  \item \textbf{Inference speed: 9.0 tok/s} on Apple MPS. Adequate for batch processing; GPU hardware or model quantization would be required for real-time interactive deployment.
\end{enumerate}

% ─────────────────────────────────────────────────────────────
\begin{thebibliography}{9}

\bibitem{hu2022lora}
Edward J.\ Hu et al.\ \textit{LoRA: Low-Rank Adaptation of Large Language Models.} ICLR 2022. \url{https://arxiv.org/abs/2106.09685}

\bibitem{dettmers2023qlora}
Tim Dettmers et al.\ \textit{QLoRA: Efficient Finetuning of Quantized LLMs.} NeurIPS 2023. \url{https://arxiv.org/abs/2305.14314}

\bibitem{rugpt3}
Sber AI. \textit{RuGPT-3 Large.} HuggingFace Model Hub. \url{https://huggingface.co/ai-forever/rugpt3large_based_on_gpt2}

\bibitem{peft}
HuggingFace. \textit{PEFT: Parameter-Efficient Fine-Tuning.} \url{https://github.com/huggingface/peft}

\bibitem{alpaca}
Rohan Taori et al.\ \textit{Stanford Alpaca: An Instruction-following LLaMA model.} 2023. \url{https://github.com/tatsu-lab/stanford_alpaca}

\end{thebibliography}

\end{document}
